\SourceLeaf{232}
\[
\begin{aligned}
{}[97]\quad&\frac1\gamma F(\alpha,\beta,\alpha+\beta+1-\gamma,1-x)\\&\qquad F(\alpha+1,\beta+1,\gamma+1,x)\\
&+\frac1{\alpha+\beta+1-\gamma}F(\alpha+1,\beta+1,\alpha+\beta+2-\gamma,1-x)\\&\qquad F(\alpha,\beta,\gamma,x)\\
&=\frac{\Pi(\alpha+\beta-\gamma)\Pi(\gamma-1)}{\Pi\alpha\Pi\beta}x^{-\gamma}(1-x)^{\gamma-\alpha-\beta-1}
\end{aligned}
\]
\SixthArticle{52}
Denotando functionem $F(-\alpha,-\beta,1-\gamma,x)$ per $S$, erit
\[
0=\alpha\beta S-(1-\gamma+(\alpha+\beta-1)x)\frac{\mathrm dS}{\mathrm dx}-(x-xx)\frac{\mathrm{dd}S}{\mathrm dx^2}
\]
Combinata hac aequatione cum I art. praec. fit
\[
0=\alpha\beta\left(\frac{S\mathrm dP+P\mathrm dS}{\mathrm dx}\right)
-(1-2x)\frac{\mathrm dP}{\mathrm dx}\mathbin{.}\frac{\mathrm dS}{\mathrm dx}
-(x-xx)\frac{\mathrm dS\mathrm{dd}P+\mathrm dP\mathrm{dd}S}{\mathrm dx^2}
\]
quae est integrabilis atque suppeditat
\[
\mathrm{Const.}=\alpha\beta PS-(x-xx)\frac{\mathrm dP}{\mathrm dx}\mathbin{.}\frac{\mathrm dS}{\mathrm dx}
\]
Valor quantitatis constantis sponte demanat $=\alpha\beta$ ex $x=0$. Si formam finitam mavis, habes
\[
\begin{aligned}
{}[98]\quad&F(\alpha,\beta,\gamma,x)F(-\alpha,-\beta,1-\gamma,x)\\
&-\frac{\alpha\beta}{\gamma-\gamma\gamma}\\&\qquad(x-xx)F(\alpha+1,\beta+1,\gamma+1,x)\\&\qquad F(1-\alpha,1-\beta,2-\gamma,x)=1
\end{aligned}
\]
Transformando hic singulas quatuor functiones secundum formulam 82 ac
dein scribendo $\gamma-\alpha,\gamma-\beta$ pro $\alpha,\beta$, habebis
\[
\begin{aligned}
{}[99]\quad&(1-x)F(\alpha,\beta,\gamma,x)F(1-\alpha,1-\beta,1-\gamma,x)\\
&-\frac{\gamma-\alpha\mathbin{.}\gamma-\beta}{\gamma-\gamma\gamma}\\&\qquad xF(\alpha,\beta,\gamma+1,x)F(1-\alpha,1-\beta,2-\gamma,x)=1
\end{aligned}
\]
\begin{center}\textit{Quaedam theoremata specialia.}\end{center}
\SixthArticle{53}
Relationes omnes quas hactenus eruimus eatenus sunt generalissimae, quod
elementa $\alpha,\beta,\gamma$ nullis circumscribuntur conditionibus. Praeterea autem plures
alias invenimus, quae conditiones speciales inter elementa $\alpha,\beta,\gamma$ supponunt:
multo vero plures sine dubio adhuc latent, easque ipsas quas hic trademus
fortasse ex altioribus principiis derivare in posterum licebit.
\par\SourceLeaf{233}
Statuamus primo in aequ. 80, $x=\frac{4y}{(1+y)^2}$ unde
\[
\mathrm dx=\mathrm dy\mathbin{.}\frac{4(1-y)}{(1+y)^3}
\]
adeoque
\[
\begin{aligned}
\frac{\mathrm dP}{\mathrm dx}&=\frac{\mathrm dP}{\mathrm dy}\mathbin{.}\frac{(1+y)^3}{4(1-y)}\\
\frac{\mathrm{dd}P}{\mathrm dx}&=\mathrm d\frac{\mathrm dP}{\mathrm dy}\mathbin{.}\frac{(1+y)^3}{4(1-y)}\\&\qquad+\frac{\mathrm dP}{\mathrm dy}\mathbin{.}\frac{(2-y)(1+y)^2}{2(1-y)^2}\mathrm dy\\
\frac{\mathrm{dd}P}{\mathrm dx^2}&=\frac{\mathrm{dd}P}{\mathrm dy^2}\mathbin{.}\frac{(1+y)^6}{16(1-y)^2}\\&\qquad+\frac{(2-y)(1+y)^5}{8(1-y)^3}\mathbin{.}\frac{\mathrm dP}{\mathrm dy}
\end{aligned}
\]
Hinc fit aequatio illa
\[
\begin{aligned}
0={}&\alpha\beta P\\
&-(\gamma(1+y)^2-4(\alpha+\beta+1)y)\frac{1+y}{4(1-y)}\mathbin{.}\frac{\mathrm dP}{\mathrm dy}\\
&-\frac{\mathrm{dd}P}{\mathrm dy^2}\mathbin{.}\frac{y(1+y)^2}4
-\frac{y(2-y)(1+y)}{2(1-y)}\mathbin{.}\frac{\mathrm dP}{\mathrm dy}
\end{aligned}
\]
sive
\[
\begin{aligned}
&0=4\alpha\beta(1-y)P\\
&-(\gamma(1+y)^2-4(\alpha+\beta+1)y-2y(2-y)(1+y))\frac{\mathrm dP}{\mathrm dy}\\
&-(y-yy)(1+y)^2\frac{\mathrm{dd}P}{\mathrm dy^2}\\
&0=4\alpha\beta(1-y)P\\
&-(1+y)(\gamma-(4\alpha+4\beta-2\gamma)y+(\gamma-2)yy)\frac{\mathrm dP}{\mathrm dy}\\
&-(1+y)^2(y-yy)\frac{\mathrm{dd}P}{\mathrm dy^2}
\end{aligned}
\]
Statuendo $P=(1+y)^{2\alpha}Q$, hinc deducitur
\[
\begin{aligned}
\mathrm I.\quad 0={}&2\alpha(2\beta-\gamma+(2\alpha+1-\gamma)y)Q\\
&-(\gamma-(4\beta-2\gamma)y+(\gamma-4\alpha-2)yy)\frac{\mathrm dQ}{\mathrm dy}\\
&-(y-yy)(1+y)\frac{\mathrm{dd}Q}{\mathrm dy^2}
\end{aligned}
\]
Iam supponendo esse $\beta=\alpha+\frac12$, haec aequatio induit formam sequentem
\[
\begin{aligned}
0={}&2\alpha(2\alpha+1-\gamma)Q\\
&-(\gamma-(4\alpha+2-\gamma)y)\frac{\mathrm dQ}{\mathrm dy}\\
&-(y-yy)\frac{\mathrm{dd}Q}{\mathrm dy^2}
\end{aligned}
\]
cuius integrale est
\[
Q=F(2\alpha,2\alpha+1-\gamma,\gamma,y)
\]
\SourceLeaf{234}
ita ut resultet
\[
\begin{aligned}&[100]\quad(1+y)^{2\alpha}F(2\alpha,2\alpha+1-\gamma,\gamma,y)\\&=F(\alpha,\alpha+\tfrac12,\gamma,\tfrac{4y}{(1+y)^2})\end{aligned}
\]
\SixthArticle{54}
Si loco relationis $\beta=\alpha+\frac12$, hanc adoptamus $\gamma=2\beta$, aequatio I art.
praec. fit
\[
\begin{aligned}
0={}&2\alpha(2\alpha+1-2\beta)yQ\\
&-(2\beta-(4\alpha+2-2\beta)yy)\frac{\mathrm dQ}{\mathrm dy}\\
&-y(1-yy)\frac{\mathrm{dd}Q}{\mathrm dy^2}
\end{aligned}
\]
Iam statuendo $yy=z$, fit
\[
\begin{aligned}
\frac{\mathrm dQ}{\mathrm dy}&=2y\frac{\mathrm dQ}{\mathrm dz}\\
\frac{\mathrm{dd}Q}{\mathrm dy^2}&=4yy\frac{\mathrm{dd}Q}{\mathrm dz^2}+\frac{2\mathrm dQ}{\mathrm dz}
\end{aligned}
\]
adeoque
\[
\begin{aligned}
0={}&\alpha(\alpha+\tfrac12-\beta)Q\\
&-(\beta+\tfrac12-(2\alpha+\tfrac32-\beta)z)\frac{\mathrm dQ}{\mathrm dz}\\
&-(z-zz)\frac{\mathrm{dd}Q}{\mathrm dz^2}
\end{aligned}
\]
cuius integrale quum sit
\[
Q=F(\alpha,\alpha+\tfrac12-\beta,\beta+\tfrac12,z)
\]
habemus
\[
\begin{aligned}&[101]\quad(1+y)^{2\alpha}F(\alpha,\alpha+\tfrac12-\beta,\beta+\tfrac12,yy)\\&=F(\alpha,\beta,2\beta,\tfrac{4y}{(1+y)^2})\end{aligned}
\]
\SixthArticle{55}
Statuamus secundo $x=4y-4yy$, unde
\[
\begin{aligned}
\mathrm dx&=4\mathrm dy(1-2y)\\
\frac{\mathrm dP}{\mathrm dx}&=\frac{\mathrm dP}{\mathrm dy}\mathbin{.}\frac1{4(1-2y)}\\
\frac{\mathrm{dd}P}{\mathrm dx^2}&=\frac{\mathrm{dd}P}{\mathrm dy^2}\mathbin{.}\frac1{16(1-2y)^2}
+\frac{\mathrm dP}{\mathrm dy}\mathbin{.}\frac1{8(1-2y)^3}
\end{aligned}
\]
unde aequatio 80 fit
\SourceLeaf{235}
\[
\begin{aligned}
&0=4\alpha\beta P\\
&-(\gamma-(4\alpha+4\beta+2)y+(4\alpha+4\beta+2)yy)\frac1{(1-2y)}\mathbin{.}\frac{\mathrm dP}{\mathrm dy}\\
&-(y-yy)\frac{\mathrm{dd}P}{\mathrm dy^2}
\end{aligned}
\]
Ut in membro secundo fractionem auferre liceat, statuere oportet $\gamma=\alpha+\beta+\frac12$,
unde prodibit
\[
\begin{aligned}
0={}&4\alpha\beta P\\
&-(\alpha+\beta+\tfrac12-(2\alpha+2\beta+1)y)\frac{\mathrm dP}{\mathrm dy}\\
&-(y-yy)\frac{\mathrm{dd}P}{\mathrm dy^2}
\end{aligned}
\]
cuius integrale est
\[
P=F(2\alpha,2\beta,\alpha+\beta+\tfrac12,y)
\]
unde habemus
\[
\begin{aligned}&[102]\quad F(\alpha,\beta,\alpha+\beta+\tfrac12,4y-4yy)\\&=F(2\alpha,2\beta,\alpha+\beta+\tfrac12,y)\end{aligned}
\]
Si in hac aequatione $y$ mutaremus in $1-y$, prodiret inde
\[
F(\alpha,\beta,\alpha+\beta+\tfrac12,4y-4yy)=F(2\alpha,2\beta,\alpha+\beta+\tfrac12,1-y)
\]
unde sequi videtur paradoxon
\[
F(2\alpha,2\beta,\alpha+\beta+\tfrac12,y)=F(2\alpha,2\beta,\alpha+\beta+\tfrac12,1-y)
\]
quae aequatio certo est falsa. Quod ut solvamus, meminisse oportet, quod probe
distinguendum est inter duas significationes characteristicae $F$, quatenus scilicet
\textit{vel} repraesentat functionem, cuius indoles exprimitur per aequationem differentialem 80, \textit{vel} solam summam seriei infinitae. Posterior, quamdiu elementum
quartum inter $-1$ et $+1$ situm est, semper exhibet quantitatem ex asse determinatam, sed cavendum est, ne hos limites excedas, quum alioquin nulla prorsus significatio supersit. Prior vero significatio repraesentat functionem generalem, quae quidem secundum legem continuitatis semper mutatur, si elementum
quartum fluxu continuo mutatur, sive ipsi valores reales sive imaginarios tribuas,
si modo semper valores 0 et 1 evites. Hinc patet, in posteriori sensu functionem
pro aequalibus elementi quarti valoribus (transitu seu potius reditu per quantitates imaginarias facto) valores inaequales adipisci posse, e quibus is quem \textit{series} $F$
\SourceLeaf{236}
repraesentat unicus tantum est, adeoque neutiquam est contradictorium, quod,
dum \textit{aliquis} valor functionis $F(\alpha,\beta,\alpha+\beta+\frac12,4y-4yy)$ est aequalis ipsi
$F(2\alpha,2\beta,\alpha+\beta+\frac12,y)$, \textit{alius} valor fit $=F(2\alpha,2\beta,\alpha+\beta+\frac12,1-y)$ foretque aeque
absurdum, hinc aequalitatem horum valorum concludere, ac si ex $\mathrm{Arc.\ sin}\frac12=30^\circ$,
$\mathrm{Arc.\ sin}\frac12=150^\circ$ concluderes $30^\circ=150^\circ$. --- Si vero characteristicam $F$ accipimus in significatione minus generali, scilicet, ut tantummodo repraesentet summam seriei $F$, ratiocinia ea, per quae aequationem 102 eruimus, necessario supponunt, $y$ tantummodo eousque a valore 0 crescere, donec evaserit $x=1$, i. e. usque ad
$y=\frac12$. In hoc ipso puncto autem \textit{continuitas} seriei $P=F(\alpha,\beta,\alpha+\beta+\frac12,4y-4yy)$
interrumpetur, quum manifesto $\frac{\mathrm dP}{\mathrm dy}$ e valore positivo (finito) subito ad negativum saliat. Itaque in hac significatione aequatio 102 extensionem ultra limites
$y=\frac12-\BareRoot{\frac12}$ usque ad $y=\frac12$ non patitur. Si mavis, eandem aequationem ita
quoque exhibere potes
\[
\begin{aligned}&[103]\quad F(\alpha,\beta,\alpha+\beta+\tfrac12,x)\\&=F(2\alpha,2\beta,\alpha+\beta+\tfrac12,\tfrac{1-\sqrt{1-x}}2)\end{aligned}
\]
sive ita
\[
\begin{aligned}&[104]\quad F(\alpha,\beta,\alpha+\beta+\tfrac12,1-x)\\&=F(2\alpha,2\beta,\alpha+\beta+\tfrac12,\tfrac{1-\BareRoot{x}}2)\end{aligned}
\]
unde tamquam corollarium sequitur (formula 48)
\[
\begin{aligned}
{}[105]\quad F(2\alpha,2\beta,\alpha+\beta+\tfrac12,\tfrac12)&=\frac{\Pi(\alpha+\beta-\tfrac12)\Pi(-\tfrac12)}{\Pi(\alpha-\tfrac12)\Pi(\beta-\tfrac12)}\\
&=\frac{\Pi(\alpha+\beta-\tfrac12)\BareRoot{\pi}}{\Pi(\alpha-\tfrac12)\Pi(\beta-\tfrac12)}
\end{aligned}
\]
\SixthArticle{56}
Ex applicatione formulae 87 ad aequationem 104 sequitur
\[
\begin{aligned}&[106]\quad F(2\alpha,2\beta,\alpha+\beta+\tfrac12,\tfrac{1-\BareRoot{x}}2)\\&=AF(\alpha,\beta,\tfrac12,x)+B\BareRoot{x}\mathbin{.}F(\alpha+\tfrac12,\beta+\tfrac12,\tfrac32,x)\end{aligned}
\]
unde patet, seriem
\[
F(2\alpha,2\beta,\alpha+\beta+\tfrac12,\tfrac{1-t}2)
\]
exhiberi posse per seriem
\[
\begin{aligned}&A+Bt+A\frac{\alpha\mathbin{.}\beta}{1\mathbin{.}\tfrac12}\mathbin{.}tt\\&+B\frac{\alpha+\tfrac12\mathbin{.}\beta+\tfrac12}{1\mathbin{.}\tfrac32}t^3\\&+A\frac{\alpha\mathbin{.}\alpha+1\mathbin{.}\beta\mathbin{.}\beta+1}{1\mathbin{.}2\mathbin{.}\tfrac12\mathbin{.}\tfrac32}t^4+\text{etc.}\end{aligned}
\]
statuendo brevitatis caussa
\SourceLeaf{237}
\[
\begin{aligned}&A=\frac{\Pi(\alpha+\beta-\tfrac12)\Pi(-\tfrac12)}{\Pi(\alpha-\tfrac12)\Pi(\beta-\tfrac12)},\\&
B=\frac{\Pi(\alpha+\beta-\tfrac12)\Pi(-\tfrac32)}{\Pi(\alpha-1)\Pi(\beta-1)}\end{aligned}
\]
Hinc colligere licet, esse
\[
\begin{aligned}&[107]\quad F(2\alpha,2\beta,\alpha+\beta+\tfrac12,\tfrac{1+\BareRoot{x}}2)\\&=AF(\alpha,\beta,\tfrac12,x)-B\BareRoot{x}\mathbin{.}F(\alpha+\tfrac12,\beta+\tfrac12,\tfrac32,x)\end{aligned}
\]
Quodsi cui haec conclusio haud satis legitima videatur, (quam tamen extra
omne dubium collocare haud difficile foret) ad eandem aequationem sequenti modo
pervenire possemus. Ex aequatione 87 fit
\[
\begin{aligned}
&F(2\alpha,2\beta,\alpha+\beta+\tfrac12,\tfrac{1+\BareRoot{x}}2)\\&=CF(2\alpha,2\beta,\alpha+\beta+\tfrac12,\tfrac{1-\BareRoot{x}}2)\\
&\quad+D\left(\frac{1-x}4\right)^{\frac12-\alpha-\beta}F(1-2\alpha,1-2\beta,\tfrac32-\alpha-\beta,\tfrac{1-\BareRoot{x}}2)
\end{aligned}
\]
statuendo brevitatis caussa
\[
\begin{aligned}&C=\frac{\Pi(\alpha+\beta-\tfrac12)\Pi(-\tfrac12-\alpha-\beta)}{\Pi(\alpha-\beta-\tfrac12)\Pi(\beta-\alpha-\tfrac12)},\\&
D=\frac{\Pi(\alpha+\beta-\tfrac12)\Pi(\alpha+\beta-\tfrac32)}{\Pi(2\alpha-1)\Pi(2\beta-1)}\end{aligned}
\]
Ex aequatione 104 autem facile deducitur
\[
\begin{aligned}
&F(1-2\alpha,1-2\beta,\tfrac32-\alpha-\beta,\tfrac{1-\BareRoot{x}}2)\\&=F(\tfrac12-\alpha,\tfrac12-\beta,\tfrac32-\alpha-\beta,1-x)\\
&=EF(\tfrac12-\alpha,\tfrac12-\beta,\tfrac12,x)\\&\qquad+G\BareRoot{x}\mathbin{.}F(1-\alpha,1-\beta,\tfrac32,x)
\end{aligned}
\]
statuendo brevitatis caussa
\[
\begin{aligned}&E=\frac{\Pi(\tfrac12-\alpha-\beta)\Pi(-\tfrac12)}{\Pi(-\alpha)\Pi(-\beta)},\\&
G=\frac{\Pi(\tfrac12-\alpha-\beta)\Pi(-\tfrac32)}{\Pi(-\tfrac12-\alpha)\Pi(-\tfrac12-\beta)}\end{aligned}
\]
Hinc rursus sequitur per aequationem 82
\[
\begin{aligned}
&F(1-2\alpha,1-2\beta,\tfrac32-\alpha-\beta,\tfrac{1-\BareRoot{x}}2)\\
&\quad=E(1-x)^{\alpha+\beta-\frac12}F(\alpha,\beta,\tfrac12,x)\\&\qquad+G\BareRoot{x}\mathbin{.}(1-x)^{\alpha+\beta-\frac12}F(\alpha+\tfrac12,\beta+\tfrac12,\tfrac32,x)
\end{aligned}
\]
His substitutis colligitur statuendo
\[
AC+DE2^{2\alpha+2\beta-1}=M,\quad BC+DG2^{2\alpha+2\beta-1}=N
\]
\[
\begin{aligned}&F(2\alpha,2\beta,\alpha+\beta+\tfrac12,\tfrac{1+\BareRoot{x}}2)\\&=MF(\alpha,\beta,\tfrac12,x)+N\BareRoot{x}\mathbin{.}F(\alpha+\tfrac12,\beta+\tfrac12,\tfrac32,x)\end{aligned}
\]
cuius \textit{forma} convenit cum aequatione 107. Iam possemus quidem e sola natura
functionis $\Pi$ derivare $M=A$, $N=-B$, quum per aequ. 55, 56 facile demonstretur, esse
\[
\begin{aligned}&C=\frac{\cos(\alpha-\beta)\pi}{\cos(\alpha+\beta)\pi},\\&
\frac{DE2^{2\alpha+2\beta-1}}A=-\frac{2\sin\alpha\pi\sin\beta\pi}{\cos(\alpha+\beta)\pi},\\&
\frac{DG2^{2\alpha+2\beta-1}}B=-\frac{2\cos\alpha\pi\cos\beta\pi}{\cos(\alpha+\beta)\pi},\end{aligned}
\]
\SourceLeaf{238}
sed hoc labore ne opus quidem est. Patet enim, statuendo $x=0$, fieri debere
\[
M=F(2\alpha,2\beta,\alpha+\beta+\tfrac12,\tfrac12)=A
\]
differentiando vero illam aequationem prodit
\[
\begin{aligned}
&x^{-\frac12}\frac{\alpha\beta}{\alpha+\beta+\tfrac12}F(2\alpha+1,2\beta+1,\alpha+\beta+\tfrac32,\tfrac{1+\BareRoot{x}}2)\\
&\quad=2\alpha\beta MF(\alpha+1,\beta+1,\tfrac32,x)\\
&\qquad+\tfrac23(\alpha+\tfrac12)(\beta+\tfrac12)N\BareRoot{x}\mathbin{.}F(\alpha+\tfrac32,\beta+\tfrac32,\tfrac52,x)\\&\qquad+\tfrac12Nx^{-\frac12}F(\alpha+\tfrac12,\beta+\tfrac12,\tfrac32,x)
\end{aligned}
\]
unde statuendo $x=0$ prodit
\[
\begin{aligned}
N&=\frac{2\alpha\beta}{\alpha+\beta+\tfrac12}F(2\alpha+1,2\beta+1,\alpha+\beta+\tfrac32,\tfrac12)\\
&=\frac{2\alpha\beta}{\alpha+\beta+\tfrac12}\mathbin{.}\frac{\Pi(\alpha+\beta+\tfrac12)\Pi(-\tfrac12)}{\Pi\alpha\Pi\beta}\\
&=\frac{\Pi(\alpha+\beta-\tfrac12)\Pi(-\tfrac32)}{\Pi(\alpha-1)\Pi(\beta-1)}=-B
\end{aligned}
\]
\SixthArticle{57}
E combinatione aequationum 106, 107 habemus itaque
\[
\begin{aligned}
{}[108]\quad&2AF(\alpha,\beta,\tfrac12,x)\\
&\quad=F(2\alpha,2\beta,\alpha+\beta+\tfrac12,\tfrac{1-\BareRoot{x}}2)\\&\qquad+F(2\alpha,2\beta,\alpha+\beta+\tfrac12,\tfrac{1+\BareRoot{x}}2)
\end{aligned}
\]
\[
\begin{aligned}
{}[109]\quad&2B\BareRoot{x}\mathbin{.}F(\alpha+\tfrac12,\beta+\tfrac12,\tfrac32,x)\\
&\quad=F(2\alpha,2\beta,\alpha+\beta+\tfrac12,\tfrac{1-\BareRoot{x}}2)\\&\qquad-F(2\alpha,2\beta,\alpha+\beta+\tfrac12,\tfrac{1+\BareRoot{x}}2)
\end{aligned}
\]
Mutando in aequatione 109 $\alpha$ in $\alpha-\frac12$, $\beta$ in $\beta-\frac12$, facile videbis, inde prodire
\[
\begin{aligned}
{}[110]\quad&\frac{\alpha-\tfrac12\mathbin{.}\beta-\tfrac12}{\alpha+\beta-\tfrac12}A\BareRoot{x}\mathbin{.}F(\alpha,\beta,\tfrac32,x)\\
&\quad=F(2\alpha-1,2\beta-1,\alpha+\beta-\tfrac12,\tfrac{1+\BareRoot{x}}2)\\&\qquad-F(2\alpha-1,2\beta-1,\alpha+\beta-\tfrac12,\tfrac{1-\BareRoot{x}}2)
\end{aligned}
\]
\clearpage\selectlanguage{german}\fancyhead[L]{\small BEMERKUNGEN.}\pdfbookmark[0]{Bemerkungen}{nachlass-bemerkungen}\SourceLeaf{239}
\begin{center}\textbf{BEMERKUNGEN.}\par\smallskip\rule{.15\linewidth}{.4pt}\end{center}
In der Handschrift bildet die \textit{Determinatio seriei nostrae per aequationem differentialem etc.} den zweiten Theil einer Abhandlung, die den Titel führt, \lq\textit{Disquisitiones generales circa functiones a serie infinita $1+\frac{\alpha.\beta}{1.\gamma}x+$ etc. pendentes auctore} \textsc{Carolo Friderico Gauss} \textit{societati regiae traditae Nov.} 1811.\rq\ Der erste Theil ist nach einer weitern Ausführung der Einzelheiten in die 1812 Jan. 30 veröffentlichte bekannte Abhandlung übergegangen. Dabei sind aus 27 Artikel 37 geworden, und die 65 numerirten Gleichungen bis zu 79 vermehrt. An jene schliessen sich die Nummern der Artikel und Gleichungen des zweiten Theils in der Handschrift an, zur Bequemlichkeit des Lesers hat man sie aber nach denen der veröffentlichten Abhandlung sich fortsetzen lassen und auch auf diese die vorkommenden Citate bezogen.

Das Handexemplar der gedruckten Abhandlung \textit{Disqu. gen. circa seriem etc.} enthält zu Art. 28 die Aufzeichnung
\[
\begin{aligned}&\textit{Die beste Definition von }\Pi\textit{ ist dass}\\&
\Pi m=\int_{-\infty}^{+\infty}e^{(m+1)x}e^{-e^x}\,\mathrm dx\end{aligned}
\]
ferner die Gleichung 58 auch noch in dieser Form
\[
\begin{aligned}&[58]\qquad \log\Pi z=(z+\tfrac12)\log z-z+\tfrac12\log2\pi\\&+\tfrac16.\tfrac12 P+\tfrac18.\tfrac23 Q+\tfrac1{10}.\tfrac34 R\\&+\tfrac1{12}.\tfrac45 S+\text{etc.}\end{aligned}
\]
die unmittelbar aus 61, 62 und 57 abgeleitet werden kann.

In einem Notizbuche ist mit Hülfe der Gleichung 58 unter der \textsc{Euler}'schen Form der Werth von $\log\Pi(10+i)$ und daraus der von $\log\Pi i$ zu
$9{,}7173075-17^\circ16'57''693\,i$ und
$\Pi i=+0{,}4980156-0{,}1549496i$ berechnet (nach dem Jahre 1847).

Der Art. 48 ist hier so wiedergegeben, wie er nach vielfachen Durchstreichungen von Worten und ganzen Sätzen in der Handschrift gelesen werden muss; nach Absicht des Verfassers dürften aber wohl noch die Worte \lq\textit{e solo theoremate binomiali}\rq\ fortzulassen sein.
\begin{flushright}\textsc{Schering.}\end{flushright}
\begin{center}\rule{.4\linewidth}{.4pt}\end{center}
